\chapter*{Lời mở đầu}
\addcontentsline{toc}{chapter}{Lời mở đầu}

Mạng truyền thông \textit{Controller Area Network} (CAN) là một thành phần quan trọng trong
các hệ thống ô tô và nhúng thời gian thực nhờ khả năng truyền dữ liệu tin cậy, độ trễ thấp
và cơ chế ưu tiên bản tin theo mã định danh. Tuy nhiên, CAN ban đầu không được thiết kế
với trọng tâm là an toàn thông tin, nên vẫn tồn tại nhiều hạn chế như thiếu xác thực nguồn
gửi, không mã hóa dữ liệu và khó ngăn chặn các hành vi phát lại, giả mạo hoặc gây ngập
lưu lượng trên bus. Khi số lượng thiết bị kết nối ngày càng tăng, các rủi ro này trở thành
một vấn đề cần được quan tâm trong quá trình nghiên cứu và triển khai hệ thống.

Từ thực tế đó, đề tài xây dựng mô hình thử nghiệm \textit{CyberCAN Defense} trên nền tảng
\textit{ESP32} kết hợp transceiver CAN nhằm mô phỏng quá trình truyền nhận bản tin, các
kịch bản tấn công cơ bản và cơ chế phát hiện bất thường trên bus. Hệ thống được phát triển
theo hướng gọn nhẹ nhưng vẫn hỗ trợ các chức năng chính như giám sát lưu lượng CAN,
mô phỏng replay, spoofing và flood, đồng thời cung cấp giao diện web để theo dõi cảnh báo
và điều khiển tập trung. Báo cáo tập trung trình bày cơ sở lý thuyết của giao thức CAN, các
vấn đề an ninh liên quan, cấu trúc phần cứng dùng ESP32 và nguyên lý hoạt động của hệ
thống đã xây dựng.
